\section{Linear momentum for a control volume}
\label{sec:momentum}

This is Newton's second law for a region the fluid flows through, and it is the most common 15-mark question in the chapter: find the force needed to hold a bend, a nozzle, a vane or a cart in place. The equation itself is one line. The marks are in the modelling: which control volume to draw, which forces act on its surface, which pressure to use, and which way each momentum flux points. A clear free-body diagram of the CV is worth more than any algebra after it.

\begin{recognise}
``Find the force on the flange bolts'', ``the force to hold the bend'', ``the anchor force'', ``the thrust'', ``the force of the jet on the plate or vane'', ``the tension in the cable holding the cart'', ``the power delivered to the moving blades''. Anything with a jet changing direction or a pipe changing direction or area, and asking for a force. Rockets and carts that accelerate because of a jet.
\end{recognise}

\begin{onsheet}
The sheet prints $\sum\vect F = \pd{}{t}\int\V\rho\,d\forall + \sum(\mdot\V)_\text{out} - \sum(\mdot\V)_\text{in}$, the jet-on-vane forms $F_x = \mdot(V_{1,x} - V_{2,x})$, $F_y = \mdot(V_{1,y} - V_{2,y})$, and $\V_r = \V_f - \V_\text{cv}$, $\mdot_\text{rel} = \rho AV_r$. It does not say that the first equation gives the force \emph{on the fluid in the CV}, while the vane forms give the force \emph{by the fluid on the vane}, which is the negative; that $\sum\vect F$ must include the pressure forces on the cut sections and gravity; that gauge pressure may be used everywhere because a uniform atmosphere exerts no net force on a closed surface; that the velocities are vectors whose components carry signs; that 1-D fluxes assume uniform velocity (otherwise multiply by $\beta$); or which mass flow a moving vane sees. None of the assumptions are printed.
\end{onsheet}

\subsection*{The model}

\paragraph{The equation.} Put $B = m\V$, so $b = \V$, into the transport theorem. For a system, Newton's law says $d(m\V)_\text{sys}/dt = \sum\vect F$, the total external force on the matter. At the instant the system fills the CV, the forces on the system are the forces on whatever is inside the CV. So, for a CV fixed in an inertial frame,
\[
\sum\vect F = \pd{}{t}\int_\text{CV}\V\rho\,d\forall + \int_\text{CS}\V\rho(\V\cdot\nhat)\,dA .
\]
This is a vector equation, three scalar equations, one per direction. With uniform velocity at each crossing it becomes
\[
\sum\vect F = \pd{}{t}\int_\text{CV}\V\rho\,d\forall + \sum_\text{out}\mdot\V - \sum_\text{in}\mdot\V .
\]
$\mdot\V$ is the \emph{momentum flux} at a section: the momentum carried across it per second. It has units $\mathrm{kg/s}\times\mathrm{m/s} = \mathrm N$, the same as a force.

\paragraph{Two separate signs.} Each term $\mdot\V$ has two signs and they are not the same thing. The sign in front, $+$ for out and $-$ for in, comes from $\V\cdot\nhat$ and says whether momentum is leaving or arriving. The sign \emph{inside}, of each component of $\V$, comes from the coordinate axes. Water entering a CV while moving in $-x$ contributes $-\mdot(-V) = +\mdot V$ to the $x$-equation. Write $\V$ in components with your axes first, then apply the out-minus-in signs. This is the single most common sign error in the chapter.

\paragraph{Which control volume.} The CV is a modelling choice, and a good one makes the unknown appear once and everything else known. Three rules. Cut where you know the pressure and velocity (at flanges with gauges, at free jets where the pressure is atmospheric). Cut \emph{through} the thing whose force you want (the bolts, the support, the cable), so that force appears as an external force on the CV. Cut perpendicular to the flow at inlets and outlets, so that $\V\cdot\nhat = \pm V$ and the pressure acts along the flow direction. For a bend held by flange bolts, the usual CV contains the bend and the water in it, with the surface cut through the bolts at each flange.

\paragraph{Which forces.} The forces on the CV are of two kinds. \emph{Body forces} act on every bit of mass inside: gravity, $\int\rho\vect g\,d\forall$, which is the weight of the fluid and of any solid inside the CV. \emph{Surface forces} act on the CS: pressure, which pushes inward along $-\nhat$ on every patch, viscous shear, which is usually negligible on short cuts, and the force from any solid the surface cuts through (a support, a bolt, a hand). At an inlet cut the upstream fluid pushes into the CV with force $P_1A_1$ along the flow; at an outlet cut the downstream fluid pushes back into the CV with $P_2A_2$, against the flow. Forgetting these pressure forces is the commonest way to get a bend force wrong by a factor of ten.

\paragraph{Why gauge pressure works.} A uniform pressure $p_a$ acting over a closed surface gives zero net force: $\oint p_a(-\nhat)\,dA = -p_a\oint\nhat\,dA = \vect 0$. The reason is geometric: for any direction, the projected areas of the parts of a closed surface facing that way and facing the opposite way are equal. (A cork in still air feels no net push from the atmosphere.) The outside of a pipe bend is surrounded by air at $p_a$, while the cut sections carry $P_1$ and $P_2$. Subtract $p_a$ everywhere: the outside surface now carries zero, and the cuts carry $P_1 - p_a$ and $P_2 - p_a$, the \emph{gauge} pressures. Since the subtracted field gave zero net force, nothing has changed. The force on the support that comes out is the real mechanical force. Use absolute pressure only if part of the CS is not surrounded by atmosphere, such as a nozzle discharging into a vacuum. A free jet in air is at atmospheric pressure, so its gauge pressure is zero.

\paragraph{Nonuniform velocity.} With a velocity profile $u$ across a section, the momentum flux is $\int\rho u^2\,dA$, not $\mdot V_\text{avg}$. Define the \emph{momentum correction factor}
\[
\beta = \frac{1}{A}\int_A\Big(\frac{u}{V_\text{avg}}\Big)^2 dA, \qquad\text{so that}\qquad \int\rho u^2\,dA = \beta\,\mdot V_\text{avg}.
\]
For laminar pipe flow $\beta = 4/3$ (derived in \secref{sec:rtt}{the transport theorem}); for turbulent pipe flow $\beta \approx 1.02$, close enough to 1 that it is usually ignored. $\beta \ge 1$ always.

\paragraph{Fixed vanes.} A jet of area $A$ and speed $V$ in air strikes a fixed curved vane and leaves at angle $\theta$ from its original direction. Along the vane the pressure is atmospheric on the free side, gravity is usually small over the short path, and friction is neglected, so Bernoulli (\secref{sec:bernoulli}{Bernoulli}) says the speed is unchanged. With $x$ along the incoming jet and the CV around the vane cut through its support,
\[
F_x = \mdot V(1 - \cos\theta),\qquad F_y = \mdot V\sin\theta
\]
for the magnitudes of the force by the fluid on the vane: $F_x$ along the incoming jet, $F_y$ away from the side toward which the jet is turned. A flat plate normal to the jet is $\theta = 90^\circ$ and gives $\mdot V$; a full reversal, $\theta = 180^\circ$, gives $2\mdot V$, the most a jet can deliver.

\paragraph{Moving vanes.} Let the vane move away from the nozzle at constant speed $u$. A frame riding on the vane is inertial, because it does not accelerate, and in that frame the flow is steady. The jet approaches at $V_r = V - u$, slides round the vane at $V_r$, and leaves at $V_r$ in the vane's direction. So the fixed-vane formulas hold with $V$ replaced by $V - u$, but with a choice of mass flow:
\begin{itemize}[nosep,leftmargin=*]
\item \emph{A single vane} running away from the nozzle sees the mass flow $\rho A(V - u)$. The jet between nozzle and vane gets longer every second, and the water filling that lengthening jet never reaches the vane. Force $\rho A(V - u)^2(1 - \cos\theta)$.
\item \emph{A wheel of many vanes} (a Pelton wheel) always has a vane in the jet, and the nozzle-to-vane distance does not grow, so every kilogram the nozzle emits is turned: mass flow $\rho AV$. Force $\rho AV(V - u)(1 - \cos\theta)$.
\end{itemize}
In both cases the power is force times vane speed, $\dot W = F_xu$. The single vane's power peaks at $u = V/3$; the wheel's at $u = V/2$, where for $\theta = 180^\circ$ it takes all the jet's kinetic energy. The absolute velocity of the leaving water is its relative velocity plus the vane velocity, $\V_2 = \V_{r,2} + \vect u$.

\paragraph{Accelerating CV.} If the CV accelerates with $\vect a_\text{rel}$ (a rocket, a cart pushed by a jet), Newton's law no longer holds in its frame. Measuring velocities relative to the CV, the law gains a fictitious force:
\[
\sum\vect F - \int_\text{CV}\vect a_\text{rel}\,\rho\,d\forall = \pd{}{t}\int_\text{CV}\V_r\rho\,d\forall + \int_\text{CS}\V_r\rho(\V_r\cdot\nhat)\,dA .
\]
The second term on the left is mass times acceleration moved to the force side. For a rocket of mass $m$ rising vertically with exhaust leaving at speed $V_e$ relative to it and exit pressure equal to ambient, this gives $m\,dV/dt = \mdot V_e - mg - \text{drag}$: the thrust is the exhaust momentum flux. Exam questions on this are short: write this equation, keep $m(t)$ if mass is lost, integrate.

\begin{firstmoves}
Draw the CV and cut through the support you want. Draw axes. List assumptions: steady (in which frame), incompressible, uniform velocity at each cut ($\beta = 1$), weight included or neglected and why, friction neglected. Use mass conservation for the unknown speed. Find unknown pressures by Bernoulli or the energy equation, never by assumption. Switch to gauge pressure and say so; a free jet is at zero gauge. Draw every force on the CV: pressure into the CV at each cut, weight, and the unknown support force with assumed positive components. Write each $\V$ in components with signs, then apply $\sum_\text{out}\mdot\V - \sum_\text{in}\mdot\V$ one direction at a time. State whether your answer is the force on the fluid, on the device, or by the support. For a moving vane, work in the vane frame and decide single vane or wheel before computing $\mdot$.
\end{firstmoves}

\subsection*{Worked examples}

\begin{example}
A $90^\circ$ reducing elbow lies in a horizontal plane. Water ($998\unit{kg/m^3}$) enters flowing in $+x$ through a $150\unit{mm}$ diameter flange at $180\unit{kPa}$ gauge and leaves flowing in $+y$ through a $100\unit{mm}$ flange into a pipe; $Q = 0.05\unit{m^3/s}$. Neglect losses in the elbow and the weight of elbow and water (it acts vertically and is carried separately). Find the horizontal force that the supports and flange bolts together must exert on the elbow.
\end{example}

\begin{proof}[Solution]
Assumptions: steady, incompressible, frictionless through the elbow, uniform velocity at each flange ($\beta = 1$), atmosphere outside the elbow, horizontal so no change of elevation.

Mass: $A_1 = \tfrac{\pi}{4}(0.15)^2 = 0.01767\unit{m^2}$, $A_2 = \tfrac{\pi}{4}(0.10)^2 = 7.854\times10^{-3}\unit{m^2}$, so $V_1 = Q/A_1 = 2.829\unit{m/s}$, $V_2 = Q/A_2 = 6.366\unit{m/s}$, and $\mdot = \rho Q = 49.9\unit{kg/s}$.

Outlet pressure: Bernoulli between the flanges (steady, incompressible, frictionless, along a streamline, same height):
\[
P_2 = P_1 + \tfrac12\rho(V_1^2 - V_2^2) = 180\,000 + \tfrac12(998)(8.006 - 40.53) = 163.8\unit{kPa}\text{ gauge}.
\]

CV: the elbow and the water inside, cut across both flanges and through the supports. Gauge pressure is used, so the outside surface carries no force. Forces on the CV: at the inlet the upstream water pushes in $+x$ with $P_1A_1 = 3181\unit{N}$; at the outlet the downstream water pushes back into the CV, in $-y$, with $P_2A_2 = 1286\unit{N}$; the supports exert the unknown $(R_x, R_y)$. Inflow velocity $(V_1, 0)$, outflow $(0, V_2)$.

$x$: $R_x + P_1A_1 = \mdot(0 - V_1)$, so $R_x = -49.9(2.829) - 3181 = -141 - 3181 = -3322\unit{N}$.

$y$: $R_y - P_2A_2 = \mdot(V_2 - 0)$, so $R_y = 49.9(6.366) + 1286 = 318 + 1286 = 1604\unit{N}$.

The supports push the elbow with $3.32\unit{kN}$ in $-x$ and $1.60\unit{kN}$ in $+y$: magnitude $3.69\unit{kN}$, at $25.8^\circ$ to the $-x$ axis toward $+y$. The water pushes the elbow equally and oppositely, in $+x$ and $-y$.

Check: the pressure terms are ten times the momentum terms, which is typical of pipe bends at moderate speed; dropping them would give $0.35\unit{kN}$, absurdly small for a pipe at $1.8$ atmospheres gauge.
\end{proof}

\begin{example}
A water jet ($998\unit{kg/m^3}$) of diameter $40\unit{mm}$ leaves a fixed nozzle at $25\unit{m/s}$ and strikes a vane mounted on a cart that moves directly away from the nozzle at a steady $10\unit{m/s}$. Relative to the vane, the water leaves at $150^\circ$ from its approach direction. Neglect friction on the vane and gravity over the vane. (a) Find the force components on the vane and the power delivered. (b) Find the absolute velocity of the water leaving the vane. (c) Repeat (a) for a wheel of vanes with the same blade speed and blade angle, and give its efficiency relative to the jet's kinetic energy flux. (d) At what blade speeds is the power largest in each case?
\end{example}

\begin{proof}[Solution]
Assumptions: steady in the vane frame, which is inertial because the cart's speed is constant; free jet at atmospheric pressure so gauge pressure is zero everywhere on the CS; frictionless, so the relative speed is unchanged along the vane; uniform jet.

(a) $A = \tfrac{\pi}{4}(0.04)^2 = 1.257\times10^{-3}\unit{m^2}$. In the vane frame the jet approaches at $V_r = 25 - 10 = 15\unit{m/s}$ and leaves at $15\unit{m/s}$ in the direction $(\cos150^\circ, \sin150^\circ) = (-0.866, 0.5)$. A single vane receives $\mdot_\text{rel} = \rho AV_r = 998(1.257\times10^{-3})(15) = 18.81\unit{kg/s}$. CV around the vane, cut through its mounting. Force on the water: $\vect F = \mdot_\text{rel}(\V_{r,2} - \V_{r,1})$, so
\[
F_x = 18.81\,(15\cos150^\circ - 15) = -526.6\unit{N},\qquad F_y = 18.81\,(15\sin150^\circ) = 141.1\unit{N}.
\]
The force by the water on the vane is the negative: $526.6\unit{N}$ in the direction of motion and $141.1\unit{N}$ in $-y$. Power $= 526.6\times10 = 5.27\unit{kW}$.

(b) $\V_2 = \V_{r,2} + \vect u = (10 - 12.99,\ 7.50) = (-2.99,\ 7.50)\unit{m/s}$, speed $8.07\unit{m/s}$, down from $25\unit{m/s}$.

(c) A wheel turns all the jet: $\mdot = \rho AV = 31.35\unit{kg/s}$. The relative velocities are unchanged, so $F_x = 31.35(15)(1 - \cos150^\circ) = 877.6\unit{N}$ and power $= 8.78\unit{kW}$. The jet carries kinetic energy at $\tfrac12\mdot V^2 = \tfrac12(31.35)(25)^2 = 9.80\unit{kW}$, so the efficiency is $8.78/9.80 = 89.6\%$.

(d) Single vane: $\dot W = \rho A(V - u)^2(1 - \cos\theta)u$; setting $d\dot W/du = 0$ gives $(V - u)(V - 3u) = 0$, so $u = V/3 = 8.33\unit{m/s}$ and $\dot W = 5.42\unit{kW}$. Wheel: $\dot W = \rho AV(V - u)(1 - \cos\theta)u$, largest at $u = V/2 = 12.5\unit{m/s}$, giving $9.14\unit{kW}$ and efficiency $(1 - \cos\theta)/2 = 93.3\%$.

Check: power at $10\unit{m/s}$ ($5.27\unit{kW}$) is just below the single-vane optimum, consistent with $10$ being close to $8.33$. The wheel can never exceed the jet's $9.80\unit{kW}$, and does not.
\end{proof}

\begin{example}
A fire nozzle is bolted to a hose of inside diameter $2.5\unit{in}$ and discharges a $1.0\unit{in}$ diameter jet into the air. The flow is $0.6\unit{ft^3/s}$ of water ($1.94\unit{slug/ft^3}$). Neglect losses in the nozzle. Find the gauge pressure at the coupling and the force in the coupling bolts, in lbf and N.
\end{example}

\begin{proof}[Solution]
Assumptions: steady, incompressible, frictionless nozzle, uniform velocity at the coupling and in the jet, nozzle horizontal, weight of nozzle and water neglected (it is vertical anyway), jet at atmospheric pressure.

Mass: $A_1 = \tfrac{\pi}{4}(2.5/12)^2 = 0.03409\unit{ft^2}$, $A_2 = \tfrac{\pi}{4}(1/12)^2 = 5.454\times10^{-3}\unit{ft^2}$. $V_1 = 0.6/0.03409 = 17.60\unit{ft/s}$, $V_2 = 0.6/0.005454 = 110.0\unit{ft/s}$, $\mdot = 1.94\times0.6 = 1.164\unit{slug/s}$.

Pressure: Bernoulli from coupling to jet, jet at zero gauge:
\[
P_1 = \tfrac12\rho(V_2^2 - V_1^2) = \tfrac12(1.94)(12\,102 - 309.8) = 11\,440\unit{lbf/ft^2} = 79.4\unit{psi}\text{ gauge}.
\]

CV: the nozzle and the water inside, cut across the coupling (through the bolts) and across the jet just outside the exit. Take $+x$ along the flow. Forces in $x$: the hose water pushes in with $P_1A_1 = 11\,440\times0.03409 = 389.9\unit{lbf}$; the jet face carries zero gauge pressure; the bolts exert $F_b$ on the nozzle. Momentum:
\[
F_b + P_1A_1 = \mdot(V_2 - V_1) = 1.164(110.0 - 17.6) = 107.6\unit{lbf},
\]
so $F_b = 107.6 - 389.9 = -282.3\unit{lbf}$. The bolts pull the nozzle upstream with $282\unit{lbf}$, so they are in tension: the water tries to push the nozzle off the hose. In SI, $282.3\times4.45 = 1256\unit{N}$.

Check: the answer is the pressure force minus the momentum gained. The water presses outward on the converging wall, whose inner face looks partly upstream, so a downstream push on the nozzle is the expected sign. $79\unit{psi}$ is a normal hydrant-fed nozzle pressure. (A slug is $\mathrm{lbf\,s^2/ft}$, so $\mathrm{slug/s}\times\mathrm{ft/s} = \mathrm{lbf}$.)
\end{proof}

\begin{trap}
Leaving out $P_1A_1$ and $P_2A_2$ at the cut sections, or giving them the wrong direction: pressure always pushes \emph{into} the CV. Using absolute pressure at the cuts while ignoring the atmosphere on the outside. Treating the sign of the momentum flux and the sign of the velocity component as one sign. Reporting the force on the fluid when the force on the device was asked. Using $\rho AV$ for a single moving vane (it is $\rho A(V - u)$), or $\rho A(V - u)$ for a wheel (it is $\rho AV$). Using absolute velocities in the vane frame. Assuming $P_2 = P_1$ in a reducer. Forgetting that the jet speed is unchanged over a frictionless vane only in the vane's frame. Writing the jet-on-vane formula as the force on the fluid with the vane's sign.
\end{trap}

\begin{ladder}
\item \white{P3.40}. A round jet strikes a fixed plate square on; the holding force. The simplest momentum balance.
\item \white{P3.64}. A jet on a plate with a central hole; the force. Two outflows with different momentum.
\item \white{P3.42}. A sudden contraction discharging to air; the force on the fitting. Pressure force at the inlet cut, in symbols.
\item \white{P3.62}. A two-branch nozzle venting angled jets; the inlet flange bolt force. Vector momentum flux with pressure at a cut.
\item \white{P3.46}. A sheet splits on a frictionless inclined plate; the split fraction. No tangential force as an assumption that closes the problem.
\item \white{P3.80}. A river over a submerged weir; the force on the obstacle. Hydrostatic pressure forces at the cuts, with mass.
\item \white{P3.88}. A pump-jet boat with drag proportional to speed squared; cruising speed. Moving control volume and a quadratic.
\item \white{P3.82}. A jet pushes a model car from rest; the speed after one metre. Accelerating control volume, integrated.
\item \white{P3.108}. A sled braked by a water scoop; the speed history. Moving vane on an accelerating body.
\end{ladder}
Answers are in the back of White; settle each call the same sitting.


\begin{hides}
Applied to an infinitesimal box, the same balance of surface forces, body force and momentum flux is the derivation of the Navier--Stokes equations in chapter 4; the pressure term there is $-\nabla P$, the gradient form of ``pressure pushes into the CV''. The momentum integral of a boundary layer, from which the sheet's K\'arm\'an relation comes, is this equation for a thin CV next to a wall. Along a single streamline the momentum equation becomes Bernoulli (\secref{sec:bernoulli}{Bernoulli}). Turbomachines use its moment form (\secref{sec:angular}{angular momentum}).
\end{hides}
